%==============================================================================
%  CAPÍTULO VIII
%==============================================================================
\chapter{Procedimiento concursal de liquidación simplificada}
\label{cap:liquidacion-simplificada}

\fichaCapitulo{Realización rápida para personas y MIPEs.}{Entrega voluntaria,
desformalización de la incautación, venta electrónica y fin de la cuenta tradicional.}

%==============================================================================
\bloqueTeorico
%==============================================================================

\subsection{Dos instituciones bajo un mismo nombre}

Bajo la denominación corriente de «liquidación simplificada» conviven dos figuras que
conviene distinguir desde el inicio, porque su ámbito y su fuente son distintos:

\begin{description}[leftmargin=0pt,style=nextline,font=\sffamily\bfseries]
  \item[La liquidación de los bienes de la Persona Deudora] Regulada en el Capítulo V de la
  ley, es el procedimiento judicial de liquidación aplicable a la \personadeudora{} y a las
  empresas deudoras que califican como micro o pequeñas. Su versión voluntaria simplificada
  se rige por el \art{273 A} y siguientes.

  \item[La realización sumaria o simplificada del activo] Regulada en los \art{203} a
  \art{206}, es un \emph{modo} de realizar los bienes ---no un procedimiento autónomo---
  que se aplica cuando el deudor califica como microempresa o cuando el producto probable
  de la realización no excede de \uf{5.000}.
\end{description}

\subsection{La humanización del concurso}

La liquidación de la persona natural planteaba, bajo el esquema tradicional, un problema de
proporcionalidad. Aplicar a una persona sobreendeudada el aparato de la quiebra mercantil
---incautación física en el hogar, martillero, cuenta detallada, espera prolongada para el
sobreseimiento--- imponía un costo humano y económico desproporcionado respecto del escaso
activo realizable.

La reforma introdujo un procedimiento aligerado que sustituye la lógica del apremio por la
de la colaboración: entrega voluntaria de bienes en lugar de incautación forzosa, subasta
electrónica en lugar de martillo presencial, y supresión de trámites que sólo tenían
sentido frente a patrimonios de magnitud.

\subsection{Del apremio a la colaboración: la entrega voluntaria de bienes}

El eje del nuevo diseño es el reemplazo de la incautación física por la entrega voluntaria.
El fundamento es doble: reduce el costo del procedimiento ---receptores, fuerza pública,
conflicto--- y reconoce que, tratándose de una persona de buena fe con escaso activo, el
apremio no aporta valor y sí produce daño.

\begin{foco}[La contrapartida de la colaboración]
La entrega voluntaria descansa sobre un presupuesto: que el deudor declaró todo. El sistema
lo incentiva reactivando la incautación forzosa cuando aparecen activos no declarados. La
humanización del procedimiento es, por tanto, un beneficio condicionado a la integridad del
inventario, y ése es el eje del trabajo forense (capítulo \ref{cap:tram-liq-simplificada}).
\end{foco}

%==============================================================================
\bloqueNormativo
%==============================================================================

\subsection{Ámbito de aplicación de la realización sumaria}

El \art{203} fija los dos supuestos en que procede la realización sumaria o simplificada:

\begin{enumerate}
  \item Si el deudor \destacar{califica como microempresa} conforme al artículo segundo de
  la \Lpyme{} \citeyearpar{ley-20416}, circunstancia que acredita el liquidador, pudiendo
  requerir al \sii{} la información sobre el nivel de ventas.
  \item Si el liquidador informa a los acreedores, en la junta constitutiva, que
  \destacar{el producto probable de la realización no excederá de \uf{5.000}}. El deudor o
  cualquier acreedor puede discrepar de la estimación.
\end{enumerate}

\subsection{Solicitud de liquidación voluntaria simplificada}

El \art{273 A} enumera los antecedentes que debe acompañar el deudor que inicia una
liquidación voluntaria simplificada. El primero es el de mayor densidad:

\begin{enumerate}
  \item \textbf{Nómina de todos los bienes de su dominio}, con indicación de su avalúo
  comercial, estado de conservación, gravámenes y ubicación, \destacar{incluyendo los que
  estén en su poder en calidad distinta de la de dueño} y los constituidos en garantía a su
  favor, con la documentación que lo acredite. Debe indicar, además, su participación en
  sociedades, comunidades y comunidades hereditarias.
  \item \textbf{Documentación de dominio} de los bienes registrables (certificado de dominio
  vigente del Conservador para inmuebles; certificado de anotaciones vigentes del Registro Civil
  para vehículos).
  \item \textbf{Nómina de bienes legalmente excluidos} del procedimiento.
  \item \textbf{Relación de juicios pendientes} con efectos patrimoniales.
  \item \textbf{Estado de deudas} ---nombre de los acreedores, naturaleza y monto de los
  créditos--- más el informe de deuda de la \cmf{}.
  \item \textbf{Nómina de trabajadores}, con prestaciones adeudadas, fueros y estado de las
  cotizaciones de seguridad social.
  \item \textbf{Cartolas históricas} de las cuentas corrientes y vistas de los dos años previos,
  emitidas dentro de los cinco días anteriores a la solicitud.
  \item \textbf{Carpeta tributaria electrónica}.
  \item \textbf{Declaración jurada} de que los antecedentes son completos y fehacientes.
\end{enumerate}

Tratándose de una \destacar{persona deudora}, la ley ordena que los antecedentes patrimoniales y
tributarios sean \destacar{reservados} ---sólo acceden a ellos el liquidador, los acreedores y la
\superir{}---, no pueden usarse con otros fines ni almacenarse, y deben eliminarse al término del
procedimiento. La declaración jurada de antecedentes fidedignos se rige por la \ncg{22}
\citeyearpar{ncg-22}. Véase también \textcite{superir-faq-liquidacion-simplificada}.

\subsection{Admisibilidad y el plazo de espera de cinco años}

El \art{273 B} establece una restricción temporal decisiva: \destacar{no puede solicitar la
liquidación voluntaria el deudor respecto del cual exista una resolución de término de una
liquidación ---ordinaria o simplificada--- firme, sino una vez transcurridos cinco años}
desde la fecha de su publicación.

El mismo artículo faculta al juez para denegar curso a la solicitud ante la insuficiencia o
el incumplimiento de los requisitos del \art{273 A}.

\begin{alerta}[El plazo de espera es de cinco años]
Una persona que obtuvo la extinción de sus saldos insolutos no puede volver a liquidarse
antes de cinco años contados desde la publicación de la resolución de término. Es un dato
que debe verificarse ---consultando procedimientos anteriores por RUT--- antes de aceptar el
encargo, porque su omisión conduce a una solicitud inadmisible. Véase el capítulo
\ref{cap:entrevista}.
\end{alerta}

\subsection{Verificación, pago y término en la liquidación de la Persona Deudora}

El Capítulo V remite, para la liquidación de los bienes de la persona deudora, a las reglas del
Capítulo IV, con una diferencia de plazo que conviene retener. El \art{277} fija la verificación
ordinaria en \destacar{quince días} contados desde la notificación de la Resolución de
Liquidación ---la mitad de los treinta días de la liquidación de la empresa deudora (\art{170})---.
La determinación del pasivo, el pago (\art{280}) y la cuenta final y término (\art{281}) se
rigen, por remisión, por las normas del Capítulo IV.

\begin{alerta}[Quince días en la persona, treinta en la empresa]
Coexisten tres plazos de verificación distintos en la ley: quince días en la reorganización
(\art{70}), treinta en la liquidación de la empresa deudora (\art{170}) y quince en la
liquidación de la persona deudora (\art{277}). Confundirlos produce verificaciones
extemporáneas. Debe identificarse siempre qué procedimiento y qué sujeto rige el caso.
\end{alerta}

\subsection{Exención de la incautación física en el hogar del deudor}

La regla general de reemplazo de la incautación por la entrega voluntaria admite
excepciones que la reactivan: la orden expresa del tribunal, la infracción de los deberes
del deudor como depositario provisional, o la detección de activos no declarados. El
tratamiento operativo de la incautación y entrega de bienes está regulado por el
Instructivo SUPERIR \Nro 5 \parencite{instructivo-i5}.

\subsection{Realización de activos mediante subasta electrónica}

El \art{204} fija las reglas de realización: los valores mobiliarios con presencia bursátil
se venden en remate en bolsa, y los demás bienes muebles e inmuebles, mediante venta al
martillo. El uso de plataformas electrónicas para la enajenación de bienes muebles se rige
por la \ncg{25} \citeyearpar{ncg-25}, que permite prescindir del martillero presencial
tradicional. Las bases de venta las confecciona el liquidador, se publican en el
\boletin{}, y el deudor o los acreedores pueden objetarlas dentro de segundo día.

\subsection{Plazos de realización y deber de información}

El \art{209} exige que la realización se efectúe en el menor tiempo posible, con un tope de
\destacar{cuatro meses para los bienes muebles} y \destacar{siete para los inmuebles},
contados desde la junta constitutiva. Los acreedores pueden acordar, con quórum calificado y
antes del vencimiento, una extensión fundada de hasta cuatro meses más.

El \art{205} impone al liquidador un deber de información: si no puede cumplir los plazos,
debe informarlo a la \superir{} con al menos quince días de anticipación al vencimiento,
justificando su demora cada treinta días, sin que ello lo exima de perseverar en la venta.
El retraso imputable al liquidador habilita las potestades sancionatorias del órgano.

\subsection{Realización forzosa de la Persona Deudora: el artículo 282}

El \art{282} regula la causal para que un acreedor solicite la liquidación de los bienes de
la persona deudora: mientras no se declare la admisibilidad de una renegociación, cualquier
acreedor puede pedirla si existen dos o más títulos ejecutivos vencidos de obligaciones
diversas, con al menos dos ejecuciones iniciadas, sin que el deudor haya presentado bienes
suficientes dentro de los cuatro días siguientes al requerimiento.

\begin{foco}[La renegociación bloquea la liquidación forzosa]
El \art{282} sólo opera «mientras no se declare la admisibilidad» de una renegociación. La
admisibilidad de la renegociación, por tanto, cierra la puerta a la liquidación forzosa que
un acreedor pretenda. Es una razón adicional para ingresar tempranamente a la renegociación
cuando el deudor califica (capítulo \ref{cap:renegociacion}).
\end{foco}

\subsection{La entrega de bienes sin incautación}

El \art{275} consagra el rasgo más humanizador del procedimiento de la persona deudora: en él
\destacar{no es necesaria la diligencia de incautación}. El liquidador se limita a requerir al
deudor la entrega de los bienes con \destacar{al menos cinco días de anticipación}. Se sustituye
así el apremio del receptor por un requerimiento de colaboración, coherente con la escasa
magnitud del activo típico y con la dignidad de la persona natural.

\begin{foco}[La entrega reemplaza a la incautación, pero la colaboración es obligatoria]
No hay incautación forzosa, pero sí deber de entrega. El deudor que se niega a entregar, o que
oculta bienes, reactiva la incautación y expone el \discharge{} al incidente de mala fe
(capítulo \ref{cap:mala-fe}). La humanización del procedimiento es un beneficio condicionado a
la buena fe.
\end{foco}

\subsection{El límite temporal del embargo de la remuneración}

El \art{276} fija una protección esencial de la persona deudora: sin perjuicio de la
inembargabilidad general del artículo 445 \Nro 2 del \cpc{}, \destacar{sólo puede embargarse la
remuneración de la persona deudora hasta por tres meses} contados desde la resolución de
liquidación de sus bienes. Si la persona deudora es casada, la realización de sus bienes se
sujeta, cuando proceda, a las normas del \cc{} sobre el régimen matrimonial.

\begin{alerta}[El sueldo no se embarga indefinidamente]
La remuneración del deudor persona natural sólo queda afecta durante tres meses desde la
resolución. Proyectar la retención más allá de ese plazo ---o incautar el sueldo íntegro sin
respetar la inembargabilidad del artículo 445 del \cpc{}--- es un error que perjudica al deudor y
compromete al liquidador. Debe verificarse, además, el régimen matrimonial, que altera qué
bienes integran la masa.
\end{alerta}

%==============================================================================
\bloquePractico
%==============================================================================

\subsection{Tratamiento de las remuneraciones del deudor persona natural}

Conforme al \textcite{instructivo-i5} ---aprobado por Resolución Exenta N.\textsuperscript{o}
14.477---, el liquidador sólo puede incautar el \destacar{exceso sobre 56 UF} de la remuneración
líquida del deudor, en línea con el límite de inembargabilidad del \artcdt{57}. La instrucción
precisa además la base de cálculo de ese exceso: los \destacar{descuentos obligatorios} por ley
(cotizaciones previsionales y de salud, impuestos) se \emph{excluyen}, mientras que los
\destacar{descuentos voluntarios} ---ahorro previsional voluntario, Cuenta 2 de la AFP, seguros,
créditos de cajas de compensación--- se \emph{incluyen}, es decir, no reducen la base incautable.

\begin{flujograma}{Tratamiento de la remuneración del deudor en la liquidación simplificada}{cap08-remuneraciones}
  \node[hito] (eval) {EVALUACIÓN DE LA REMUNERACIÓN DEL DEUDOR};

  \coordinate (split) at ($(eval.south)+(0,-6mm)$);
  \node[favorable, below left=12mm and 8mm of eval.south] (hasta)
    {HASTA \uf{56}\\[2pt]\normalfont Inembargable absoluto: prohibición de incautar};
  \node[adverso, below right=12mm and 8mm of eval.south] (mas)
    {MÁS DE \uf{56}\\[2pt]\normalfont Tramo excedente sujeto a incautación};

  \node[nodo, text width=34mm, below=9mm of hasta] (hastadesc)
    {Descuentos obligatorios y voluntarios: excluidos de la base};
  \node[nodo, text width=34mm, below=9mm of mas] (masdesc)
    {Obligatorios excluidos; voluntarios incluidos en la base del excedente};

  \draw[flecha] (eval.south) -- (split) -| (hasta.north);
  \draw[flecha] (eval.south) -- (split) -| (mas.north);
  \draw[flecha] (hasta) -- (hastadesc);
  \draw[flecha] (mas) -- (masdesc);
\end{flujograma}

\begin{practica}[Controversias con el empleador]
Si surge discrepancia interpretativa entre el liquidador y el empleador respecto de la
retención de los excedentes, el asunto debe someterse al tribunal del concurso por la vía
incidental del \art{131} ---audiencias verbales a solicitud del interesado, que expone por escrito
su petición y antecedentes--- , mecanismo general al que remiten las controversias sobre la
administración de los bienes, incluida la incautación de remuneraciones (capítulo
\ref{cap:liquidacion-ordinaria}).
\end{practica}

\subsection{Los bienes inembargables: el límite de la masa incautable}

La incautación de la persona deudora encuentra un límite infranqueable en el catálogo de bienes
\destacar{inembargables} del \artcpc{445}, que el liquidador debe respetar bajo su responsabilidad.
Los de mayor incidencia en el concurso de la persona natural son:

\begin{description}[leftmargin=0pt,style=nextline,font=\sffamily\bfseries]
  \item[Pensiones (N.\textsuperscript{o} 1)] Sueldos, gratificaciones y \destacar{pensiones de
  gracia, jubilación, retiro y montepío} que pagan el Estado y las municipalidades ---embargables
  sólo hasta el 50\,\% por deudas de pensiones alimenticias---.
  \item[Remuneraciones (N.\textsuperscript{o} 2)] En la forma de los \artcdt{40} y \artcdt{153},
  concordante con el límite de las 56 UF (capítulo \ref{cap:liquidacion-ordinaria}).
  \item[Alimentos forzosos (N.\textsuperscript{o} 3)] Las pensiones alimenticias que el deudor
  recibe.
  \item[La vivienda familiar (N.\textsuperscript{o} 8)] El bien raíz que el deudor ocupa con su
  familia \destacar{si su avalúo fiscal no supera 50 UTM}, más el mobiliario de dormitorio, comedor
  y cocina y la ropa de uso familiar. La inembargabilidad del inmueble no rige, con todo, frente al
  Fisco, las cajas de previsión y los organismos del \textsc{minvu}.
  \item[Instrumentos de trabajo (N.\textsuperscript{os} 9 a 12)] Libros de la profesión, máquinas
  de enseñanza, uniformes militares y útiles del arte u oficio ---hasta 50 UTM y a elección del
  deudor---, y los aperos y animales de labor del trabajador de campo.
\end{description}

\begin{alerta}[Incautar un bien inembargable expone al liquidador]
El catálogo del \artcpc{445} no es una recomendación: incautar la pensión de jubilación, la
vivienda de avalúo inferior a 50 UTM o las herramientas de trabajo del deudor es una actuación
antijurídica que da lugar a la tercería, a la restitución y a la responsabilidad del liquidador. La
depuración de lo incautable ---especialmente en la persona deudora, cuyo patrimonio suele reducirse
a bienes exentos--- debe hacerse \emph{antes} de la diligencia, no después del reclamo.
\end{alerta}

%==============================================================================
\bloqueJurisprudencial
%==============================================================================

\subsection{Denegación del \discharge{} por omisión de activos realizables}

Sentencias que facultan al tribunal para declarar un \discharge{} parcial o denegarlo
íntegramente cuando el deudor oculta cuentas de ahorro previsional voluntario o no
declara vehículos motorizados vigentes, por infracción de la buena fe
objetiva.\verificar{Individualizar los fallos y precisar si la denegación parcial tiene
sustento en el texto vigente o es construcción pretoriana.}
